\section*{अध्याय \ref{ch:g5:natural-hazards} --- प्राकृतिक संकट}

\begin{solution}{exo:g5:natural-hazards:1}
(क) केवल \omterm{def:g5:natural-hazards:hazard}{प्राकृतिक संकट}: किसी को नुक़सान नहीं होता। (ख) \omterm{def:g5:natural-hazards:disaster}{प्राकृतिक आपदा}: घरों को
क्षति पहुँचती है। (ग) केवल \omterm{def:g5:natural-hazards:hazard}{प्राकृतिक संकट}। (घ) \omterm{def:g5:natural-hazards:disaster}{प्राकृतिक आपदा}।
\end{solution}

\begin{solution}{exo:g5:natural-hazards:2}
हवा से: लू, \omterm{def:g5:natural-hazards:storm}{तूफ़ान}, \omterm{def:g5:natural-hazards:drought}{सूखा}। पानी से: \omterm{def:g5:water-shapes-the-land:flood}{बाढ़}। ज़मीन से: \omterm{def:g5:natural-hazards:landslide}{भूस्खलन}, \omterm{def:g4:volcanoes-earthquakes-first-look:volcano}{ज्वालामुखी} \omterm{def:g4:volcanoes-earthquakes-first-look:eruption}{उद्गार},
\omterm{def:g4:volcanoes-earthquakes-first-look:earthquake}{भूकंप}।
\end{solution}

\begin{solution}{exo:g5:natural-hazards:3}
\omterm{def:g5:natural-hazards:drought}{सूखा} महीनों या वर्षों तक का एक लंबा समय है, जिसमें सामान्य से बहुत कम \omterm{def:g4:clouds-rain-snow:precipitation}{बारिश} होती
है। वह किसानों को बिना फ़सल के छोड़ सकता है, क्योंकि फ़सलें मुरझाकर मर जाती हैं, और
शहरों में पीने के पानी की कमी कर सकता है, क्योंकि \omterm{def:g1:landscapes-around-us:water}{नदियाँ} और \omterm{def:g1:landscapes-around-us:water}{झीलें}
सिकुड़ जाती हैं।
\end{solution}

\begin{solution}{exo:g5:natural-hazards:4}
उदाहरण के लिए: जो कुछ भी हवा उड़ा ले जा सकती है, उसे अंदर ले आना या बाँध देना
(बग़ीचे की कुर्सियाँ, खिलौने, कूड़ेदान); टॉर्च जाँचना और अतिरिक्त बैटरियाँ और पानी
तैयार रखना; \omterm{def:g5:natural-hazards:storm}{तूफ़ान} के दौरान घर के अंदर, खिड़कियों से दूर रहना; ख़बरों के लिए रेडियो
सुनते रहना।
\end{solution}

\begin{solution}{exo:g5:natural-hazards:5}
(क) हरिकेन। (ख) टाइफ़ून। (ग) उष्णकटिबंधीय चक्रवात।
\end{solution}

\begin{solution}{exo:g5:natural-hazards:6}
हरा, पीला, नारंगी, लाल। लाल स्तर पर बड़ा ख़तरा है: तुरंत क़दम उठाओ और ज़िम्मेदार
लोगों के निर्देश मानो।
\end{solution}

\begin{solution}{exo:g5:natural-hazards:7}
दाईं ओर की \omterm{def:g1:landscapes-around-us:landform}{पहाड़ी} की ढलान का एक हिस्सा टूटकर नीचे फिसल गया, और पीछे एक नंगा घाव
छोड़ गया। उसका \omterm{def:g3:stones-pebbles-sand:sand}{कीचड़} \omterm{def:g1:landscapes-around-us:landform}{घाटी} की पूरी तली पर फैल गया और सड़क को ढक दिया। उसने \omterm{def:g1:landscapes-around-us:water}{नदी} को
भी एक बाँध की तरह रोक दिया: \omterm{def:g1:landscapes-around-us:water}{नदी} का पानी, जो अब आगे नहीं बह सकता, उसके पीछे जमा हो
गया है और उसने ऊपर की ओर के जंगल में \omterm{def:g5:water-shapes-the-land:flood}{बाढ़} ला दी है।
\end{solution}

\begin{solution}{exo:g5:natural-hazards:8}
लीना सही है। हम \omterm{def:g4:clouds-rain-snow:precipitation}{बारिश} या \omterm{def:g1:landscapes-around-us:water}{नदी} को नहीं रोक सकते, पर लोगों को ख़तरे से दूर रख सकते
हैं। उदाहरण के लिए: घर \omterm{def:g5:water-shapes-the-land:flood}{बाढ़} के \omterm{def:g1:landscapes-around-us:landform}{मैदान} के बजाय ऊँची ज़मीन पर बनाना; \omterm{def:g5:water-shapes-the-land:flood}{बाढ़} की
चेतावनियों से लोगों को समय पर चेताना, ताकि वे ऊँची ज़मीन पर चले जाएँ; सबको
सिखाना कि \omterm{def:g5:water-shapes-the-land:flood}{बाढ़} के पानी में से कभी पैदल न जाएँ और न गाड़ी
चलाएँ।
\end{solution}

\begin{solution}{exo:g5:natural-hazards:9}
(क) $4 \times 60 = 240$ मिलीमीटर। (ख) 60 का एक-चौथाई 15 है; चार
महीनों में $4 \times 15 = 60$ मिलीमीटर। (ग) $240 - 60 = 180$ मिलीमीटर
कम पड़े: सामान्य \omterm{def:g4:clouds-rain-snow:precipitation}{बारिश} का तीन-चौथाई।
\end{solution}

\begin{solution}{exo:g5:natural-hazards:10}
(क) $1\,200 \div 4 = 300$ घर। (ख) 300 का एक-तिहाई 100 है; इसलिए
$300 - 100 = 200$ घर अब भी \omterm{def:g5:water-shapes-the-land:flood}{बाढ़} के \omterm{def:g1:landscapes-around-us:landform}{मैदान} पर हैं।
\end{solution}

\begin{solution}{exo:g5:natural-hazards:11}
संभव कारण: पहले गाँव को समय पर चेताया गया था और लोगों ने तैयारी कर ली थी
(खिड़कियों के पल्ले बंद, चीज़ें अंदर), दूसरे को नहीं; पहले गाँव के घर मज़बूती से बने
हैं, दूसरे के कमज़ोर; दूसरा गाँव किसी नीचे \omterm{def:g1:landscapes-around-us:water}{तट} या \omterm{def:g5:water-shapes-the-land:flood}{बाढ़} के \omterm{def:g1:landscapes-around-us:landform}{मैदान} पर है, जहाँ \omterm{def:g5:natural-hazards:storm}{तूफ़ान}
\omterm{def:g5:water-shapes-the-land:flood}{बाढ़} भी लाया, जबकि पहला ऊँची ज़मीन पर है; दूसरे गाँव में घरों के पास के पेड़ उन पर
गिर पड़े। एक ही संकट हर जगह एक-सा नुक़सान नहीं करता: यह इस पर निर्भर है कि लोग कहाँ
और कैसे रहते हैं, और उन्होंने कितनी अच्छी
तैयारी की है।
\end{solution}
